% Spacetime Emergence: 7D Gamma -> 3+1D M^4
% -----------------------------------------------------------------
% How the seven dimensions project onto physical spacetime.
%
% The SU(3) decomposition 7 = 3 + 3_bar + 1 groups the dimensions as:
%   Space    {A, S, D}  -> fundamental 3       -> Sigma^3 (T-119)
%   Internal {L, E, U}  -> conjugate 3_bar     -> SM gauge fibre
%   Time     {O}        -> singlet 1           -> R (T-118)
%
% To make the three groups contiguous in the top row, dimensions are
% arranged in visual order (A, S, D, L, E, U, O) -- NOT the Fano-
% canonical order (A, S, D, L, E, O, U). This is purely a layout
% convenience: dimension labels themselves are unchanged.
% -----------------------------------------------------------------

\begin{tikzpicture}[
  dimbox/.style={rectangle, rounded corners=3pt,
                 minimum width=1.2cm, minimum height=0.8cm,
                 align=center, font=\small\bfseries},
  grp/.style={rounded corners=6pt, dashed, thick},
  arr/.style={-{Stealth[length=3mm]}, very thick, #1},
  scale=1.0
]

% === Title ===
\node[font=\large\bfseries] at (0, 5.6)
  {$\Gamma \in \mathcal{D}(\mathbb{C}^7)$: seven dimensions};

% === 7D row, reordered as {A,S,D} | {L,E,U} | {O} for contiguous grouping ===
% Centres at x = -4.5, -3.0, -1.5, 0, 1.5, 3.0, 4.5
% Box width 1.2 -> half-width 0.6, so edges:
%   A  : [-5.1, -3.9]
%   S  : [-3.6, -2.4]
%   D  : [-2.1, -0.9]
%   L  : [-0.6,  0.6]
%   E  : [ 0.9,  2.1]
%   U  : [ 2.4,  3.6]
%   O  : [ 3.9,  5.1]
\node[dimbox, draw=blue,           fill=blue!15]   (A) at (-4.5, 4) {A};
\node[dimbox, draw=blue,           fill=blue!15]   (S) at (-3.0, 4) {S};
\node[dimbox, draw=blue,           fill=blue!15]   (D) at (-1.5, 4) {D};
\node[dimbox, draw=orange,         fill=orange!15] (L) at ( 0.0, 4) {L};
\node[dimbox, draw=red,            fill=red!15]    (E) at ( 1.5, 4) {E};
\node[dimbox, draw=purple,         fill=purple!15] (U) at ( 3.0, 4) {U};
\node[dimbox, draw=green!50!black, fill=green!15]  (O) at ( 4.5, 4) {O};

% === Grouping dashed boxes (non-overlapping, with 0.1 gaps) ===
% All three rectangles span y in [3.50, 4.50] so every dim box
% (which sits at y = [3.6, 4.4]) is vertically centred in its group.
%
% {A,S,D} group:  x in [-5.20, -0.80]  (contains A,S,D at [-5.1, -0.9])
% {L,E,U} group:  x in [-0.70,  3.70]  (contains L,E,U at [-0.6,  3.6])
% {O}     group:  x in [ 3.80,  5.20]  (contains O     at [ 3.9,  5.1])
% Gaps of 0.10 between consecutive groups, no overlap anywhere.
\draw[blue,           grp] (-5.20, 3.50) rectangle (-0.80, 4.50);
\draw[orange,         grp] (-0.70, 3.50) rectangle ( 3.70, 4.50);
\draw[green!50!black, grp] ( 3.80, 3.50) rectangle ( 5.20, 4.50);

% === Group labels (compact, centred just below each group box) ===
\node[font=\scriptsize, blue, text width=3.8cm, align=center,
      inner sep=2pt]
  at (-3.0, 2.75)
  {\textbf{Space} $\{A,S,D\}$\\
   fundamental $\mathbf{3}$ of $\mathrm{SU}(3)$};

\node[font=\scriptsize, orange, text width=3.8cm, align=center,
      inner sep=2pt]
  at ( 1.5, 2.75)
  {\textbf{Internal} $\{L,E,U\}$\\
   $\bar{\mathbf{3}}$: gauge\,+\,exp.\,+\,unity};

\node[font=\scriptsize, green!50!black, text width=1.3cm, align=center,
      inner sep=2pt]
  at ( 4.5, 2.75)
  {\textbf{Time} $\{O\}$\\
   singlet $\mathbf{1}$};

% === Projection arrows with tight midway labels ===
% Labels are anchored to the arrow midpoint with near-zero
% horizontal gap (inner sep = 2pt), eliminating the wide gap of
% the previous version.
\draw[arr=blue!60] (-3.0, 1.95) -- (-3.0, -0.05)
  node[midway, left=2pt, font=\footnotesize, blue!60,
       text width=2.1cm, align=right]
  {Gelfand\,+\,Connes\\ reconstr.\ (T-119)};

\draw[arr=green!50!black] (4.5, 1.95) -- (4.5, -0.05)
  node[midway, right=2pt, font=\footnotesize, green!50!black,
       text width=2.1cm, align=left]
  {Page--Wootters\\ clock algebra (T-118)};

% === M^4 level ===
\node[rectangle, rounded corners=8pt, draw=plospurple, fill=plospurple!10,
      minimum width=11cm, minimum height=1.3cm, align=center, font=\bfseries]
  (M4) at (0, -0.75)
  {$M^{4} = \mathbb{R}\times\Sigma^{3}$ \quad (T-120)};

\node[font=\footnotesize, color=plosgray] at (0, -1.85)
  {Product spectral triple $\to$ Einstein--Hilbert action
   $+$ Standard Model (T-121)};

% === Vacuum topology ===
\node[font=\footnotesize, color=plospurple] at (0, -2.45)
  {Vacuum: $\Sigma^{3}\cong S^{3}$ (de Sitter), $\Lambda > 0$
   \quad (T-120b)};

% === Key insight (bottom caption inside figure) ===
\node[font=\footnotesize\itshape, color=gray,
      text width=11cm, align=center]
  at (0, -3.30) {
  Spacetime is not fundamental. It is the macroscopic projection of
  the 7-dimensional coherence matrix onto the
  $\{A,S,D\}\!\cup\!\{O\}$ sector. Gravity is coherence seen from the
  external aspect.
};

\end{tikzpicture}
